\documentclass[10pt]{article}

\usepackage{parskip}
\usepackage[margin=1in]{geometry} 
\usepackage{tikz}
\usetikzlibrary{positioning}
\usepackage{amsmath,amsthm,amssymb, graphicx, multicol, array}
\usepackage{enumitem}
\usepackage{amssymb}
\usepackage{float}
\usepackage{href-ul}
 
\newcommand{\N}{\mathbb{N}}
\newcommand{\Z}{\mathbb{Z}}
 
\newenvironment{problem}[2][Problem]{\begin{trivlist}
\item[\hskip \labelsep {\bfseries #1}\hskip \labelsep {\bfseries #2.}]}{\end{trivlist}}

\begin{document}
 
\title{Experimental Design}
\author{Jakob Sverre Alexandersen\\
GRA4158 Marketing and the Analysis of Experiments and Quasi-Experiments\\
Lecture 6}
\maketitle

\tableofcontents
\newpage
\section{Manipulating the independent variable}
\begin{itemize}
    \item How do we vary the treatment? I.e., the independent variable 
    \begin{itemize}
        \item Selecting the right treatment and control condition(s)
    \end{itemize}
    \item What is the counterfactual?
    \begin{itemize}
        \item Is it the absence of something?
        \item Or is it an explicit alternative?
        \item What is the focal concept that we want to test / research?
        \item How can we correctly implement the alternatives without inference?
    \end{itemize}
    \item Balance between realism and a clear definition of the desired manipulation
    \begin{itemize}
        \item Internal vs external validity
    \end{itemize}
    \item Variation of the independent variable must meet two requirements: 
    \begin{itemize}
        \item Effectiveness: the treatment must be effective in the sense that the focal concept is indeed varied 
        \item Control: the treatment must gurantee that any aspects other than the focal independent variable are not unintentionally varied 
    \end{itemize}
\end{itemize}
\section{Examples}
\subsection{Scarcity cues in E-commerce}
Causal question: does displaying scarcity influence purchase intentions?
\begin{itemize}
    \item IV: Scarcity
    \begin{itemize}
        \item Treatment: deisplaying in-stock information
        \item counterfactual/control: ??
    \end{itemize}
    \item Effectiveness: is displaying in-stock information actually varying scarcity perception?
    \begin{itemize}
        \item Is the displayed number large/small enough to affect scarcity in a meaningful way?
        \item Manupulation check: measure scarcity?
    \end{itemize}
    \item control: is displaying in-stock information varying any other variables than scarcity? like popularity etc. 
\end{itemize}
\subsection{Compound treatment --- example}
\begin{itemize}
    \item you want to investigate the causal effect of a price discount on purchase decisions 
    \item you run an A/B test 
    \begin{itemize}
        \item You randomly assign customers to two groups 
        \item customers in the treatment group get an email with the price discount 
        \item customers in the control group do not get an email 
    \end{itemize}
    \item \textbf{is the difference between groups the effect of the price discount?}
\end{itemize}
\newpage
\subsection{Similar: placebo effects}
\begin{itemize}
    \item You want to test the effect of receiving a medicine for pain relief 
    \item you run an A/B test 
    \begin{itemize}
        \item you randomly assign customers to two groups 
        \item people in the treatment get the pill
        \item people in the control group do not get the pill
    \end{itemize}
    \item \textbf{is the difference between groups the effect of the medicine?}
    \item no, because observed differences between the treatment and control group may be due to receiving the pull and not because the medicine actually works 
\end{itemize}
\section{Measuring the dependent variable(s)}
\begin{itemize}
    \item how can we find a way to measure the dependent variable(s)? the intended outcome variable(s)
    \item what is the variable that we are interested in?
    \begin{itemize}
        \item are we able to measure it direcly? Do we need proxy variables to measure it indirectly?
        \item are we able to measure it in all treatment (and control) conditions in the same way? 
        \item is measurement possible in the absence of treatment?
    \end{itemize}
    \item does the type of measure introduce threats to validity? demand effects, social desiribility etc. 
    \item do we need to capture additional variables to understand the process and mechanism of the treatment better (i.e., mediators)?
\end{itemize}
\section{validity of experiments}
\begin{minipage}[t][0.28\textheight][t]{0.48\textwidth}
\textbf{Internal validity}
\begin{itemize}
    \item refers to whether the treatment actually caused the observed effects on the dependent variable(s)
    \item can we accurately setimate the treatment effect within the given context and for our particular example?
    \begin{itemize}
        \item fails when there are differences between treatment groups (other than what we can control) that affect the outcome
        \item directly associated with the design of the experiment
    \end{itemize}
    \vfill
\end{itemize}
\end{minipage}
\hfill
\begin{minipage}[t][0.28\textheight][t]{0.48\textwidth}
\textbf{External validity}
\begin{itemize}
    \item refers to whether the cause-and-effect relationships found in the experiment can be generalized 
    \item can we extrapolate our estimates to other settings/populations/time?
    \begin{itemize}
        \item fails when the treatment effect is different
        \begin{itemize}
            \item outside of the evalutation environment 
            \item or for other groups of people
        \end{itemize}
        \item how realistic / representative is the experiment for real environments in which the treatment should be applied?
    \end{itemize}
    \vfill
\end{itemize}
\end{minipage}
\vspace{1em}
\newpage 
\subsection{internal validity}
\subsubsection{Some common threats to internal validity}
\begin{itemize}
    \item The experimental protocol
    \begin{itemize}
        \item potential outcomes not independent of assignment: presence of selection effects 
        \item non-compliance, attrition 
        \item interference / violation of SUTVA
        \item time effects 
    \end{itemize}
    \item awareness about the experiment (hawthorne effect)
    \item the measurement of the outcome 
    \begin{itemize}
        \item differential measures 
        \item ambiguous temporal precedence 
    \end{itemize}
    \item statistical conclusion validity
\end{itemize}
\subsubsection{Threats to internal validity --- the experimental protocol}
\begin{itemize}
    \item potential outcomes not independent / selection effects 
    \begin{itemize}
        \item between-subjects design 
        \begin{itemize}
            \item assignment not (fully) random 
            \begin{itemize}
                \item we can test on observed pre-treatment variables for equivalence of groups 
                \item we may need to rewuire conditional independence (i.e., after controlling for observables we achieve indepenence)
            \end{itemize}
            \item non-compliance
        \end{itemize}
        \item within-subject design 
        \begin{itemize}
            \item time effects: order and carryover effects, changes in the environment etc. 
        \end{itemize}
    \end{itemize}
    \item violation of SUTVA (stable unit value assumption)
    \begin{itemize}
        \item interference: potential outcomes are not independent of the treatment status of other individuals 
    \end{itemize}
    \item \textbf{differntial attrition}
    \begin{itemize}
        \item attrition: participants leave the experiment before all measures are completed 
        \item differntial attrition: drop out in group is larger or smaller than in the other group due to reasons of the experiment / drop out is not independent of the potential outcomes 
    \end{itemize}
    \item \textbf{time effects}
    \begin{itemize}
        \item when treatments are administered at two different times, outside events, learning or other changes in the environment can influence the treatment effect
    \end{itemize}
\end{itemize}
\newpage 
\subsubsection{Stratified random assignment}
\begin{itemize}
    \item does not randomize on the population level, but within more homogenous subgroups
    \begin{itemize}
        \item assigns loyalty card members with .5 probability to treatment and control
        \item assigns non-loyalty card members with .5 probability to treatment and control
    \end{itemize}
    \item ensures better balance for each of the characteristics on which we build subgroups / strata 
    \begin{itemize}
        \item characteristics should be potentially important confounders where balancing is important
    \end{itemize}
    \item important if a key characteristic has low prevalence in the population and small sample
    \begin{itemize}
        \item example: we have a sample of 100 customers of an airline of which 2 are heavy frequent flyers. We split the population into two groups. therefore it is unlikely that the two flyers are in different groups 
    \end{itemize}
    \item esp. important if we use this characteristic to analyze conditional treatment effects 
\end{itemize}
\subsubsection{Threats to internal validity --- awareness about the experiment}
\textbf{Hawthorne effect}
\begin{itemize}
    \item if individuals know that they are part of an experiment, they might behave differently
\end{itemize}
\textbf{Demand effects}
\begin{itemize}
    \item participants guess the purpose or goald of the experiment and try to conform with it 
    \item they change their behavior in favor or against the hypothesis
\end{itemize}

Ideally people should be unaware of the experiment or the goals of the experiment, but consider ethical consequences 
\subsubsection{Threats to internal validity --- measurement of the outcome}
\textbf{Differential measures}
\begin{itemize}
    \item sometimes it is hard to measure the same outcome in different treatment groups
    \item psychological measures may change meaning due to treatment (or absence thereof)
    \item actual behavior might be harder to observe in the absence or presence of the treatment
\end{itemize}
\textbf{ambiguous temporal precedence}
\begin{itemize}
    \item the type of measurement does not allow to distinguish whether treatment actually lies before the outcome or after $\to $ resolution of the measurement is too broad
\end{itemize}
\subsubsection{Threats to internal validity --- statistical conclusion validity}
are the statistical assumptions of the method in use fulfilled?
\begin{itemize}
    \item knowledge of the design is important to choose the right statistical framwork for estimating the causal effects under given circumstances 
\end{itemize}
\newpage 
\subsection{External validity}
\subsubsection{Most common threats to external validity of experiments}
\begin{itemize}
    \item Non-representative sample
    \begin{itemize}
        \item sample size too small 
        \item use of special participant groups: students, volunteers, pre-testers etc 
    \end{itemize}
    \item non-representative design / environment
    \begin{itemize}
        \item the treatment differs in actual implementations
        \item context differs / idiosyncracy of the test situation 
        \item different outcome variables 
        \item equilibrium effects: implementing the treatment might change the economic environment as consumers and competitors adapt to the new situation 
    \end{itemize}
\end{itemize}
\subsubsection{Types of experiments}
\begin{minipage}[t][0.28\textheight][t]{0.48\textwidth}
    \textbf{Lab experiment}
    \begin{itemize}
        \item environment is created by the researcher and thus differs from reality
        \item high control over the experimental environment
        \begin{itemize}
            \item often high internal validity
        \end{itemize}
        \item low realism, participants may feel observed or do not represent the target population 
        \vfill
    \end{itemize}
    \end{minipage}
    \hfill
    \begin{minipage}[t][0.28\textheight][t]{0.48\textwidth}
    \textbf{Field experiment}
    \begin{itemize}
        \item everyday life situation; the same setting where the findings from the experiment will be deployed 
        \item less control over the experimental environment
        \begin{itemize}
            \item lower internal validity
        \end{itemize}
        \item high realism, participants usually not aware, and more representative of the target population
        \vfill
    \end{itemize}
    \end{minipage}
    \vspace{1em}
\subsubsection{Four key features of field experiments}
\begin{itemize}
    \item real-world context 
    \item authenticity of treatment
    \item representativeness of participants
    \item relevant outcome measures 
\end{itemize}
\newpage 
\section{Within subject designs}
\begin{itemize}
    \item builds on the core idea of temporal precedence
    \item so far, 
\end{itemize}
\end{document}